\documentclass[9pt,a4paper]{extarticle}

\usepackage[utf8]{inputenc}
\usepackage[T1]{fontenc}
\usepackage{mathpazo} % Font yang nyaman dibaca
\usepackage[margin=1.5cm]{geometry} % Margin lega
\usepackage{amsmath,amssymb,amsfonts}
\usepackage{multicol}
\usepackage{enumitem}
\usepackage{titlesec}
\usepackage{array}
\usepackage{xcolor} % Untuk warna
\usepackage[most]{tcolorbox} % Untuk kotak cantik dan anti-pecah antar kolom

% --- PENGATURAN SPASI ---
\setlength{\parindent}{0pt}
\setlength{\parskip}{4pt} % Jarak antar paragraf lebih lega
\linespread{1.1} % Spasi baris dilonggarkan
\setlist[itemize]{noitemsep, topsep=2pt, leftmargin=1.3em, label=\textbullet}
\setlist[enumerate]{noitemsep, topsep=2pt, leftmargin=1.5em}

% --- PENGATURAN JUDUL ---
\titleformat{\section}
  {\normalfont\large\bfseries\color{blue!70!black}}{\thesection}{0.5em}{}
  [\vspace{-4pt}\color{blue!70!black}\rule{\linewidth}{0.8pt}\vspace{4pt}]
\titleformat{\subsection}
  {\normalfont\normalsize\bfseries\color{black!80}}{\thesubsection}{0.5em}{}

% --- KOTAK RUMUS (Pakai tcolorbox biar halus di mata) ---
\newtcolorbox{rumusbox}{
    colback=blue!5!white,
    colframe=blue!30!white,
    arc=2mm,
    boxrule=0.5pt,
    left=3pt, right=3pt, top=3pt, bottom=3pt,
    fontupper=\small,
    breakable % Biar kotak bisa pindah halaman
}

\begin{document}

\begin{center}
	\textbf{\Large CHEATSHEET ATS MATEMATIKA DISKRIT} \\
	\vspace{2pt}
	\small Pertemuan 1--7 | Open Book | Nama: ................ | NIM: ................
\end{center}
\vspace{-4mm}
\noindent\rule{\linewidth}{1pt}
\vspace{2mm}

\begin{multicols}{2}

% ==========================================
\section*{1. TEORI HIMPUNAN}
Himpunan adalah kumpulan objek berbeda yang terdefinisi jelas (seperti Array / Tabel Database).
\begin{itemize}
    \item \textbf{Kardinalitas (\(|A|\)):} Jumlah elemen unik. \texttt{Contoh: } $A=\{2,3,5,7\}$, maka $|A|=4$.
    \item \textbf{Subset (\(A \subseteq B\)):} Semua anggota A ada di B.
    \item \textbf{Power Set (\(P(A)\)):} Himpunan yang isinya semua subset dari A. Rumus jumlahnya: \(2^{|A|}\).
    \item \textbf{Himpunan Kosong (\(\varnothing\)):} Himpunan tanpa elemen, \(|\varnothing|=0\).
\end{itemize}

\textbf{Operasi Himpunan:}
\begin{itemize}
    \item \textbf{Irisan (\(\cap\)):} `INNER JOIN` (DAN).
    \item \textbf{Gabungan (\(\cup\)):} `UNION` (ATAU).
    \item \textbf{Selisih (\(-\)):} `EXCEPT` (A yang bukan B).
    \item \textbf{Beda Simetris (\(\oplus\)):} Elemen di salah satu, tapi tidak keduanya.
    \item \textbf{Kartesian (\(\times\)):} Pasangan berurutan \((a,b)\). Jumlahnya $|A| \times |B|$.
\end{itemize}

\begin{rumusbox}
\textbf{\textcolor{blue!70!black}{Prinsip Inklusi-Eksklusi (Wajib Hafal!)}} \\
Menghitung gabungan tanpa menghitung ganda. \\
\textbf{2 Himpunan:} \\
\(|A \cup B| = |A| + |B| - |A \cap B|\) \\
\textbf{3 Himpunan:} \\
\(|A \cup B \cup C| = |A|+|B|+|C| - |A \cap B| - |A \cap C| - |B \cap C| + |A \cap B \cap C|\)
\end{rumusbox}

\textbf{Himpunan Fuzzy:} Berbeda dengan himpunan tegas (Crisp: 0 atau 1), Fuzzy punya derajat keanggotaan \(\mu \in [0,1]\). Contoh: Usia 40 tahun \(\mu=0.5\) (setengah muda, setengah tua).

% ==========================================
\section*{2. MATRIKS}
Susunan angka dalam baris dan kolom. Ordo = ukuran (Baris \(\times\) Kolom).
\begin{itemize}
    \item \textbf{Penjumlahan/Pengurangan:} Ordo (ukuran) harus sama persis.
    \item \textbf{Perkalian Skalar (\(k \cdot A\)):} Semua angka di dalam matriks dikali \(k\).
    \item \textbf{Perkalian Matriks (\(A \times B\)):} Syarat mutlak: \textbf{Kolom A = Baris B}. (Baris dikali Kolom).
    \item \textbf{Transpose (\(A^T\)):} Tukar posisi baris dan kolom.
\end{itemize}

\begin{rumusbox}
\textbf{\textcolor{blue!70!black}{Determinan \& Invers (Matriks Persegi)}} \\
\textbf{Det 2x2:} \(\begin{vmatrix} a & b \\ c & d \end{vmatrix} = ad - bc\) \\
\textbf{Invers 2x2:} \(A^{-1} = \frac{1}{ad-bc} \begin{pmatrix} d & -b \\ -c & a \end{pmatrix}\) \\
\textbf{Sifat Invers:} \((AB)^{-1} = B^{-1} A^{-1}\) (urutan dibalik!). \\
\textbf{Syarat Invers:} Determinan \(\neq 0\). \\
\textbf{Persamaan:} Jika \(AX = B\), maka \(X = A^{-1} B\).
\end{rumusbox}

% ==========================================
\section*{3. RELASI}
Aturan yang memasangkan anggota himpunan A ke himpunan B (Subset dari perkalian kartesian). Bisa direpresentasikan dengan Diagram Panah, Tabel, atau Matriks Biner (isi 1 jika terhubung, 0 jika tidak).

\begin{rumusbox}
\textbf{\textcolor{blue!70!black}{4 Sifat Relasi (Cek pakai Matriks)}} \\
\begin{itemize}
    \item \textbf{Refleksif:} Semua elemen menunjuk dirinya sendiri. (Diagonal matriks isinya 1 semua).
    \item \textbf{Simetris:} Jika A \(\to\) B, maka B \(\to\) A. (Matriks Biner = Transposenya).
    \item \textbf{Tolak-Setangkup:} Jika A \(\to\) B dan B \(\to\) A, maka A dan B adalah orang yang sama.
    \item \textbf{Transitif:} Jika A \(\to\) B, dan B \(\to\) C, maka harus ada A \(\to\) C.
\end{itemize}
\end{rumusbox}

\textbf{Invers Relasi (\(R^{-1}\)):} Membalik arah panah \((a,b)\) jadi \((b,a)\). \\
\textbf{Komposisi (\(S \circ R\)):} Jalankan R dulu, baru S. Dibaca dari \textbf{Kanan ke Kiri}.

% ==========================================
\section*{4. FUNGSI}
Relasi "Satu-Satu". Artinya, \textbf{setiap input di Domain hanya boleh punya SATU output} di Kodomain. (Kalau 1 input punya 2 output, itu bukan fungsi).

\begin{rumusbox}
\textbf{\textcolor{blue!70!black}{3 Sifat Fungsi (Penting untuk Invers):}}
\begin{itemize}
    \item \textbf{Injektif (One-to-One):} Output tidak ada yang kembar.
    \item \textbf{Surjektif (On-to):} Semua target di Kodomain kebagian pasangan.
    \item \textbf{Bijektif:} Keduanya terpenuhi. \textbf{Hanya fungsi Bijektif yang punya Invers (\(f^{-1}\))}.
\end{itemize}
\end{rumusbox}

\textbf{Fungsi Khusus:}
\begin{itemize}
    \item \textbf{Lantai \(\lfloor x \rfloor\):} Bulatkan ke bawah. \(\lfloor 3.7 \rfloor = 3\).
    \item \textbf{Langit \(\lceil x \rceil\):} Bulatkan ke atas. \(\lceil 3.2 \rceil = 4\).
    \item \textbf{Modulo \(a \bmod n\):} Sisa pembagian. \(17 \bmod 5 = 2\).
    \item \textbf{Rekursif:} Fungsi memanggil dirinya sendiri. Harus punya \textbf{Basis} (kondisi berhenti) dan \textbf{Rekurens} (rumusnya).
\end{itemize}

% ==========================================
\section*{5. LOGIKA MATEMATIKA}
Fondasi komputer. Komputer hanya mengerti Benar (1) dan Salah (0).

\textbf{5 Operator Wajib:}
\begin{itemize}
    \item \textbf{Negasi (\(\neg\)):} Kebalikan. (Kode: \texttt{!A})
    \item \textbf{Konjungsi (\(\wedge\)):} DAN. Benar jika keduanya Benar. (Kode: \texttt{\&\&})
    \item \textbf{Disjungsi (\(\vee\)):} ATAU. Salah jika keduanya Salah. (Kode: \texttt{||})
    \item \textbf{Implikasi (\(\to\)):} JIKA-MAKA. Hanya Salah jika (Benar \(\to\) Salah). (Kode: \texttt{!A || B})
    \item \textbf{Biimplikasi (\(\leftrightarrow\)):} JIKA DAN HANYA JIKA. Benar jika nilai A dan B sama. (Kode: \texttt{A == B})
\end{itemize}

\begin{rumusbox}
\textbf{\textcolor{blue!70!black}{Hukum Ekivalensi (Penyederhanaan):}}
\begin{itemize}
    \item \textbf{De Morgan:} \(\neg(A \wedge B) \equiv \neg A \vee \neg B\) \\
    \textit{Tukar DAN/ATAU, lalu negasi semuanya.}
    \item \textbf{Implikasi:} \(A \to B \equiv \neg A \vee B\).
    \item \textbf{Distributif:} \(A \wedge (B \vee C) \equiv (A \wedge B) \vee (A \wedge C)\).
\end{itemize}
\end{rumusbox}

\subsection*{Tipe Ekspresi \& Inferensi}
\begin{itemize}
    \item \textbf{Tautologi:} Selalu Benar (semua baris tabel B).
    \item \textbf{Kontradiksi:} Selalu Salah (semua baris tabel S).
    \item \textbf{Kontingensi:} Campuran B dan S.
    \item \textbf{Modus Ponens:} \(A \to B\) (benar), \(A\) (benar) \(\implies\) Kesimpulan: \(B\).
    \item \textbf{Modus Tollens:} \(A \to B\) (benar), \(\neg B\) (benar) \(\implies\) Kesimpulan: \(\neg A\).
\end{itemize}

\vspace{3mm}
\begin{center}
    \textit{*Fokus baca kotak warna biru saat ujian. Semoga sukses ATS-nya!*}
\end{center}

\end{multicols}
\end{document}